\chapter{肌肉驱动生物力学建模与实现}
\section{建模总体设计}
本系统建模工作基于本地项目工程环境，依托 User-in-the-Box 开源框架与 MuJoCo 动力学仿真引擎开展。建模过程充分复用框架自带的bm\_models生物力学模型库，结合本毕业设计方向盘操控任务的专属需求，对人体上肢模型、交互对象模型进行针对性修改与适配，最终实现人体肌肉与方向盘机构的稳定物理耦合，为上层控制算法、视觉引导与人机交互测试提供真实可靠的仿真载体。

相较于传统纯运动学模型，本次搭建的肌肉骨骼模型完整还原人体上肢肌肉发力、关节联动、力传递的生物力学特性，能够真实模拟肌肉驱动的自然操控手感，契合课题肌肉驱动生物力学人机交互的核心定位。
\section{系统模块与建模职责}
为便于开发与维护，系统建模按功能划分为四大模块，各模块职责如下表所示。
\begin{table}[htbp]
  \centering
  \caption{建模模块及主要职责}
  \label{tab:module_duty}
  \begin{tabular}{ll}
    \toprule
    模块名称 & 主要职责 \\
    \midrule
    数据处理模块 & 运动数据加载、清洗、标准化、对齐、增强等预处理 \\
    运动学模块 & 基于 SMPL 模型前向运动学计算，输出关节位置、速度等状态 \\
    控制模块 & 构建 Actor-Critic 架构，基于 PPO 生成关节控制信号，优化策略 \\
    仿真模块 & 基于 MuJoCo 搭建肌肉骨骼系统，执行控制并更新物理状态 \\
    评估模块 & 从轨迹相似度、生物合理性、任务完成度评估控制效果 \\
    \bottomrule
  \end{tabular}
\end{table}

\section{人体上肢肌肉骨骼模型构建}
人体上肢肌肉骨骼模型来自 User-in-the-Box 框架里的 bm\_models 生物力学模型文件夹，模型用 MuJoCo 可以识别的 XML 格式进行结构化编写，依靠 User-in-the-Box 开源框架完成模型调用和仿真环境的搭建，不需要从头制作骨骼结构，很大程度降低了人体生物力学建模的难度，本研究选用 User-in-the-Box 框架自带的人体上肢骨骼模型，借助 MuJoCo 仿真引擎实现模型加载和画面显示，如图~\ref{fig:upper_limb_model}所示，模型完整包含胸廓、脊柱和两侧上肢的骨骼结构，还设置了符合人体活动范围的关节限制条件，能够模拟肩部、肘部、腕部的多方向运动，模型通过连接 VR 手柄设备，模拟出双手握住方向盘的初始姿势，为后续肌肉驱动控制和转向交互仿真提供了贴近真实的生物力学条件。
\begin{figure}[htbp]
  \centering
  \includegraphics[width=0.8\textwidth]{image_7.png}
  \caption{Simulato.xml模型}
  \label{fig:upper_limb_model}
\end{figure}  

模型文件被放在指定的 bm\_models 文件夹里，文件夹中包含模型主体的 XML 文件、外观材质设置、碰撞几何参数和肌肉布局定义等配套资源，各个文件之间有规范的关联关系，能够被 MuJoCo 仿真引擎直接读取并加载，在仿真搭建的过程中，系统通过路径读取的方式导入这个目录下的模型文件，完成模型坐标校准、关节限位设置和物理参数初始化，实现人体上肢模型在仿真场景里的定位和静态载入。

模型导入完成以后，研究人员对上肢模型的关节自由度、肌肉发力增益、运动阻尼等重要参数进行适配修改，让模型更适合方向盘旋转操控的任务场景，同时结合 MuJoCo 的动力学计算能力，完成模型重力、惯性、接触碰撞等物理属性的配置，保证上肢模型在跟随方向盘转动时，动作顺畅自然，不会出现姿态变形和关节超出限制的问题，为后续的肌肉驱动控制、视觉引导跟踪和精度测试提供稳定可靠的模型条件。
\section{肌肉–方向盘耦合驱动实现}
本研究设计的肌肉与方向盘耦合驱动系统，是依靠 User-in-the-Box 框架和 MuJoCo 仿真引擎来实现的，系统以 mobl\_arms\_bimanual\_motor 双臂上肢模型作为主要载体，搭建出肌肉激活信号到上肢关节力矩再到方向盘运动状态的完整对应关系，为后续的方向盘转向控制和精度测试提供了稳定好用的仿真环境；耦合驱动的实现过程会以生物力学模型作为基础，mobl\_arms\_bimanual\_motor 模型在标准骨骼结构上整合了上肢主要肌肉的发力部分和关节力矩计算端口，完整还原了人体胸廓、脊柱以及两边肩膀、手肘、手腕关节的身体结构和运动限制，能比较真实地模仿上肢肌肉在操作方向盘时一起发力的过程；研究还通过逆动力学方式把肌肉激活度信号转成关节驱动的力矩，让生物力学方面的驱动基础得到实现\cite{li2026numerical}；在此基础之上，研究借助仿真场景文件 simulation.xml 完成了上肢模型和方向盘的物理连接，把方向盘模型和上肢手部关节进行固定连接，让方向盘的转动自由度和上肢手腕关节、肩关节的转动自由度形成物理约束，同时对模型初始握住方向盘的姿势进行调整，让物理参数相互匹配，保证上肢转动和方向盘转向可以同步进行，避免运动时出现姿势变形或者关节超出范围的情况；耦合驱动整个实现过程可以分成模型加载与初始化、肌肉信号对应、力矩输出与联动控制、实时仿真与验证这四个阶段，研究使用 Python 代码切换到项目根目录来加载场景文件，保证骨骼、肌肉和方向盘的物理参数设置准确，再依靠 effort\_models.py 模块把方向盘目标转向角度转成上肢关节的目标姿势，计算出肌肉激活度信号并对应成关节驱动力矩，带动上肢关节转动，同时带着方向盘一起转向，再通过 MuJoCo 引擎的实时动力学计算对耦合系统的运动状态进行更新；仿真运行界面如图~\ref{fig:coupling_simulation}所示，界面中间可以看到双臂上肢骨骼模型和方向盘握把的连接状态，界面右边的关节面板里列出了所有可以控制的自由度，界面下方的性能曲线可以实时显示仿真求解器的迭代收敛情况和 CPU 占用情况，这些内容都验证了耦合驱动系统的实时性和稳定性，也为后续基于肌肉信号的方向盘转向控制算法设计和精度测试提供了完整的物理环境和数据支持。
\begin{figure}[htbp]
  \centering
  \includegraphics[width=0.8\textwidth]{image_8.png}
  \caption{肌肉–方向盘耦合驱动仿真界面}
  \label{fig:coupling_simulation}
\end{figure}  

\section{模型合理性验证}
为确保肌肉与方向盘耦合驱动模型的物理效果和仿真表现符合要求，本研究从结构合理程度、动力变化特征、动作限制条件三个方面开展相关验证工作，研究人员会使用 MuJoCo 仿真工具自带的模型检查功能，对上肢模型和方向盘模型进行完整检测，骨骼连接关系、关节限制条件、模型固定方式都会被逐一核对，重复设置或连接中断的问题会被及时排除，模型用到的 STL 文件与 XML 配置文件也会被对应检查，文件路径错误、代码格式错误、参数设置错误都会被处理，模型可以被正常加载，整体基础结构也能保持稳定状态\cite{tu2026jiyu}。

在动力学特征验证环节，无外力作用下的自由下落与静止仿真会被用于测试，模型在重力作用下的运行状态与静止时的稳定效果会被重点观察，上肢模型在初始姿态下不会发生位置穿透、偏移等问题，方向盘与手部的连接结构也不会出现异常晃动或脱离现象，关节力矩加载实验会被用于测试，肌肉控制信号到关节动力输出的对应关系会被验证，不同信号强度下的力矩输出曲线会被对比分析，力矩输出过程平稳无突变，符合生物力学相关理论规律，肌肉驱动的运行逻辑也被证实是有效可行的。

在运动学限制验证环节，肩部、肘部、腕部关节的全范围动作测试会被开展，关节限制范围与动作空间的合理程度会被确认，上肢模型跟随方向盘转动时不会出现关节超范围运动或姿态变形问题，连续转向仿真测试会被用于记录数据，方向盘转动角度与上肢关节角度的同步变化情况会被统计，两者动作平稳无滞后，耦合连接关系稳定，与真实驾驶时肢体和方向盘的配合状态保持一致，综合各项验证结果可知，本模型在结构完整、动力响应、动作限制方面均具备良好合理性，可以满足后续肌肉驱动转向控制与精度测试的仿真需要。

